\documentclass{amsart}
% For dealing with references we use the comment environment
\usepackage{verbatim}
\newenvironment{reference}{\comment}{\endcomment}
%\newenvironment{reference}{}{}
\newenvironment{slogan}{\comment}{\endcomment}
\newenvironment{history}{\comment}{\endcomment}

% For commutative diagrams we use Xy-pic
\usepackage[all]{xy}

% We use 2cell for 2-commutative diagrams.
\xyoption{2cell}
\UseAllTwocells

% We use multicol for the list of chapters between chapters
\usepackage{multicol}

% This is generally recommended for better output
\usepackage{lmodern}
\usepackage[T1]{fontenc}

% For cross-file-references

% Package for hypertext links:
\usepackage{hyperref}

% For any local file, say "hello.tex" you want to link to please

% Theorem environments.
%
\theoremstyle{plain}
\newtheorem{theorem}[subsection]{Theorem}
\newtheorem{proposition}[subsection]{Proposition}
\newtheorem{lemma}[subsection]{Lemma}

\theoremstyle{definition}
\newtheorem{definition}[subsection]{Definition}
\newtheorem{example}[subsection]{Example}
\newtheorem{exercise}[subsection]{Exercise}
\newtheorem{situation}[subsection]{Situation}

\theoremstyle{remark}
\newtheorem{remark}[subsection]{Remark}
\newtheorem{remarks}[subsection]{Remarks}

\numberwithin{equation}{subsection}

% Macros
%
\def\lim{\mathop{\mathrm{lim}}\nolimits}
\def\colim{\mathop{\mathrm{colim}}\nolimits}
\def\Spec{\mathop{\mathrm{Spec}}}
\def\Hom{\mathop{\mathrm{Hom}}\nolimits}
\def\Ext{\mathop{\mathrm{Ext}}\nolimits}
\def\SheafHom{\mathop{\mathcal{H}\!\mathit{om}}\nolimits}
\def\SheafExt{\mathop{\mathcal{E}\!\mathit{xt}}\nolimits}
\def\Sch{\mathit{Sch}}
\def\Mor{\mathop{\mathrm{Mor}}\nolimits}
\def\Ob{\mathop{\mathrm{Ob}}\nolimits}
\def\Sh{\mathop{\mathit{Sh}}\nolimits}
\def\NL{\mathop{N\!L}\nolimits}
\def\CH{\mathop{\mathrm{CH}}\nolimits}
\def\proetale{{pro\text{-}\acute{e}tale}}
\def\etale{{\acute{e}tale}}
\def\QCoh{\mathit{QCoh}}
\def\Ker{\mathop{\mathrm{Ker}}}
\def\Im{\mathop{\mathrm{Im}}}
\def\Coker{\mathop{\mathrm{Coker}}}
\def\Coim{\mathop{\mathrm{Coim}}}

% Boxtimes
%
\DeclareMathSymbol{\boxtimes}{\mathbin}{AMSa}{"02}

%
% Macros for moduli stacks/spaces
%
\def\QCohstack{\mathcal{QC}\!\mathit{oh}}
\def\Cohstack{\mathcal{C}\!\mathit{oh}}
\def\Spacesstack{\mathcal{S}\!\mathit{paces}}
\def\Quotfunctor{\mathrm{Quot}}
\def\Hilbfunctor{\mathrm{Hilb}}
\def\Curvesstack{\mathcal{C}\!\mathit{urves}}
\def\Polarizedstack{\mathcal{P}\!\mathit{olarized}}
\def\Complexesstack{\mathcal{C}\!\mathit{omplexes}}
% \Pic is the operator that assigns to X its picard group, usage \Pic(X)
% \Picardstack_{X/B} denotes the Picard stack of X over B
% \Picardfunctor_{X/B} denotes the Picard functor of X over B
\def\Pic{\mathop{\mathrm{Pic}}\nolimits}
\def\Picardstack{\mathcal{P}\!\mathit{ic}}
\def\Picardfunctor{\mathrm{Pic}}
\def\Deformationcategory{\mathcal{D}\!\mathit{ef}}

\setlength{\emergencystretch}{3em}
\begin{document}
\title[Stacks Project, Tag 0ASJ]{426 --- Etale local extensions of valuation rings are weakly unramified}
\maketitle
\noindent Copyright (C) 2005--2025 Johan de Jong. Stacks Project authors. Source: \href{https://stacks.math.columbia.edu/tag/0ASJ}{Tag 0ASJ}.

\noindent Editorial difficulty note (not author text). Difficulty H2: After proving valuation-ring structure from weak dimension, the proof must show equality of value groups. It approximates the base by normal Noetherian local subrings and descends roots up to units through codimension-one unramified extensions, recovering the value-group statement..

\noindent Editorial provenance note (not author text). History: excerpt from revision \texttt{a0444\allowbreak 6e57e\allowbreak c1fbc\allowbreak 252a8\allowbreak 71afc\allowbreak ec775\allowbreak 2fb28\allowbreak 07b14}, more-algebra.tex, lines 37479--37532. Reference presentation was adapted; original-author context and core support blocks were retained or added. The mathematical wording is unchanged. Licensed under GNU FDL 1.2 or later; no invariant sections or cover texts. See ../licenses/COPYING. Context blocks reproduce author definitions and supporting results; they are not additional counted problems.

\begin{lemma}
\label{lemma-etale-extension-valuation-ring}
Let $A$ be a valuation ring. Let $A \to B$ be an \'etale ring map
and let $\mathfrak m \subset B$ be a prime lying over the maximal
ideal of $A$. Then $A \subset B_\mathfrak m$ is an extension of
valuation rings which is weakly unramified.
\end{lemma}

\begin{proof}
The ring $A$ has weak dimension $\leq 1$ by
Lemma \href{https://stacks.math.columbia.edu/tag/092S}{lemma weak dimension at most 1 (Tag 092S)}. Then $B$ has weak dimension
$\leq 1$ by Lemmas \href{https://stacks.math.columbia.edu/tag/092E}{lemma weak dimension goes up (Tag 092E)} and
\href{https://stacks.math.columbia.edu/tag/092N}{lemma when weakly etale (Tag 092N)}. hence the local ring $B_\mathfrak m$
is a valuation ring by Lemma \href{https://stacks.math.columbia.edu/tag/092S}{lemma weak dimension at most 1 (Tag 092S)}.
Since the extension $A \subset B_\mathfrak m$ induces a finite
extension of fraction fields,
we see that the $\Gamma_A$ has finite index in the value group of
$B_{\mathfrak m}$. Thus for every $h \in B_\mathfrak m$ there exists
an $n > 0$, an element $f \in A$, and a unit $w \in B_\mathfrak m$
such that $f = w h^n$ in $B_\mathfrak m$. We will show that this implies
$f = ug^n$ for some $g \in A$ and unit $u \in A$; this will show that
the value groups of $A$ and $B_\mathfrak m$ agree, as claimed in the lemma.

\medskip\noindent
Write $A = \colim A_i$ as the colimit of its local subrings which
are essentially of finite type over $\mathbf{Z}$. Since $A$
is a normal domain (Algebra, Lemma \href{https://stacks.math.columbia.edu/tag/00IC}{lemma valuation ring normal (Tag 00IC)}),
we may assume that each $A_i$ is normal (here we use that taking
normalizations the local rings remain essentially of finite type
over $\mathbf{Z}$ by
Algebra, Proposition \href{https://stacks.math.columbia.edu/tag/0335}{proposition ubiquity nagata (Tag 0335)}).
For some $i$ we can find an \'etale extension $A_i \to B_i$
such that $B = A \otimes_{A_i} B_i$, see
Algebra, Lemma \href{https://stacks.math.columbia.edu/tag/00U2}{lemma etale (Tag 00U2)}.
Let $\mathfrak m_i$ be the intersection of $B_i$ with $\mathfrak m$.
Then we may apply Lemma \href{https://stacks.math.columbia.edu/tag/0ASI}{lemma extension normal domains and roots (Tag 0ASI)}
to the ring map $A_i \to (B_i)_{\mathfrak m_i}$ to conclude.
The hypotheses of the lemma are satisfied because:
\begin{enumerate}
\item $A_i$ and $(B_i)_{\mathfrak m_i}$ are Noetherian as they are
essentially of finite type over $\mathbf{Z}$,
\item $A_i \to (B_i)_{\mathfrak m_i}$ is flat as $A_i \to B_i$ is \'etale,
\item $B_i$ is normal as $A_i \to B_i$ is \'etale, see
Algebra, Lemma \href{https://stacks.math.columbia.edu/tag/033C}{lemma normal goes up (Tag 033C)},
\item for every height $1$ prime of $A_i$ there exists a height $1$
prime of $(B_i)_{\mathfrak m_i}$ lying over it by
Algebra, Lemma \href{https://stacks.math.columbia.edu/tag/02MA}{lemma finite in codim 1 (Tag 02MA)} and the fact that
$\Spec((B_i)_{\mathfrak m_i}) \to \Spec(A_i)$ is surjective,
\item the induced extensions $(A_i)_\mathfrak p \to (B_i)_\mathfrak q$
are unramified for every prime $\mathfrak q$ lying over a prime
$\mathfrak p$ as $A_i \to B_i$ is \'etale.
\end{enumerate}
This concludes the proof of the lemma.
\end{proof}
\end{document}
